\documentclass[10pt]{article}
\usepackage[T1]{fontenc}
\usepackage[utf8]{inputenc}
\usepackage[french]{babel}
\usepackage{amsmath,amssymb,mathrsfs,mathpazo}
\usepackage{geometry}
\geometry{a4paper,margin=2cm}
\usepackage{graphicx,array,tabularx,tabularray,colortbl,multirow,multicol,booktabs,enumitem,pifont,ulem,fancybox,diagbox,cancel}
\usepackage{pstricks,pstricks-add,pst-plot,pst-tree,pst-node,pst-text,pst-func,pst-math,pst-bspline,pst-3dplot,pst-eucl,pst-solides3d}
\let\pscalculate\undefined
\usepackage{pst-bezier}
\usepackage{tikz,pgfplots,tkz-tab}
\usetikzlibrary{arrows,arrows.meta,patterns,calc,positioning,trees,decorations.markings,shapes.geometric,shapes.misc,shapes.arrows,intersections,decorations.pathmorphing,decorations.pathreplacing,backgrounds,shadows,babel}
\pgfplotsset{compat=1.18}
\usepgfplotslibrary{fillbetween}
\usepackage[np]{numprint}
\usepackage{listings,makecell,pas-tableur,twoopt}
\let\barre\undefined
\usepackage[pstricks]{tablvar}
\definecolor{darkgreen}{RGB}{0,100,0}
\providecommand{\red}{\color{red}}
\providecommand{\blue}{\color{blue}}
\providecommand{\green}{\color{green}}
\providecommand{\black}{\color{black}}
\providecommand{\white}{\color{white}}
\providecommand{\gray}{\color{gray}}
\providecommand{\cyan}{\color{cyan}}
\providecommand{\magenta}{\color{magenta}}
\providecommand{\yellow}{\color{yellow}}
\providecommand{\violet}{\color{violet}}
\providecommand{\orange}{\color{orange}}
\providecommand{\euro}{\texteuro}
\providecommand{\EUR}{\texteuro}
\providecommand{\ieme}{\textsuperscript{e}}
\providecommand{\degres}{\ensuremath{^\circ}}
\providecommand{\degree}{\ensuremath{^\circ}}
\providecommand{\vv}[1]{\overrightarrow{#1}}
\providecommand{\Coord}[1]{#1}
\providecommand{\decosix}{}
\providecommand{\starredbullet}{\ensuremath{\bullet}}
\pagestyle{empty}
\setlength{\parindent}{0pt}
\usepackage{adjustbox}\DeclareUnicodeCharacter{2236}{:}\DeclareUnicodeCharacter{2190}{\ensuremath{\leftarrow}}
\providecommand{\R}{\mathbb{R}}
\providecommand{\N}{\mathbb{N}}
\providecommand{\Z}{\mathbb{Z}}
\providecommand{\Q}{\mathbb{Q}}
\providecommand{\D}{\mathbb{D}}
\providecommand{\C}{\mathbb{C}}
\providecommand{\e}{{\mathrm{e}}}
\providecommand{\vect}[1]{\overrightarrow{#1}}
\providecommand{\vectt}[1]{\overrightarrow{\mathrm{#1}}}
\providecommand{\vvt}[1]{\overrightarrow{\mathrm{#1}}}
\providecommand{\np}[2][]{{#2~#1}}
\providecommand{\vv}[1]{{\overrightarrow{#1}}}
\renewcommand{\d}{{\mathrm{d}}}
\providecommand{\degres}{{^\circ}}
\providecommand{\degre}{{^\circ}}
\providecommand{\blue}{\color{blue}}
\providecommand{\red}{\color{red}}
\providecommand{\enskip}{\quad}
\providecommand{\llbracket}{\lbrack\!\lbrack}
\providecommand{\rrbracket}{\rbrack\!\rbrack}
\providecommand{\euro}{{\text{€}}}
\providecommand{\no}{{\text{n°}}}
\providecommand{\interval}[2]{{[#1;#2]}}
\providecommand{\intervFO}[2]{{[#1;#2[}}
\providecommand{\dotfill}{\ldots}

\providecommand{\ds}{\displaystyle}
\providecommand{\textding}[1]{#1}
\providecommand{\Oij}{\(\mathrm{O};\vec{\imath},\vec{\jmath}\)}
\providecommand{\Oijk}{\(\mathrm{O};\vec{\imath},\vec{\jmath},\vec{k}\)}
\providecommand{\Ouv}{\(\mathrm{O};\vec{u},\vec{v}\)}
\providecommand{\intervFF}[2]{\ensuremath{[#1;#2]}}
\renewcommand{\euro}{\ifmmode\text{\texteuro}\else\texteuro\fi}
\DeclareUnicodeCharacter{00B0}{\ifmmode^{\circ}\else\textdegree\fi}
\DeclareUnicodeCharacter{20AC}{\ifmmode\text{\texteuro}\else\texteuro\fi}
% Énoncés : Éducation nationale. Transcription et corrigés : APMEP, Denis Vergès et contributeurs. Découpage et classement : CoopMaths, Sébastien Lozano, Éric Elter et contributeurs. Adaptation de la présentation : Maths Explorer.
% CC BY-SA : https://coopmaths.fr/www/about/
% https://coopmaths.fr/alea/static/bac/2024/tex/bac_2024_05_sujet2_ameriquenord_1_cor.tex
\begin{document}
\noindent\textit{APMEP / CoopMaths. CC BY-SA. Présentation adaptée par Maths Explorer.}\par\medskip
\textbf{Partie A}

\medskip











\begin{enumerate}
\item \phantom{21}  
\begin{center}
\pstree[treemode=R,levelsep = 3.5cm]{\TR{}}
{\pstree{\TR{$N~~$}\taput{\np{0,2286}}}
	{\TR{$R$}\taput{\np{0,0808}}
	\TR{$\overline{R}$}\tbput{\red 0,9192}
	}
\pstree{\TR{$\overline{N}~~$}\tbput{\red\np{0,7714}}}
	{\TR{$R$}\taput{\np{0,0127}}
	\TR{$\overline{R}$}\tbput{\red 0,9873}
	}
}
\end{center}

\item 
On calcule $p(N \cap R) = p(N) \times p_N(R) = \np{0,2286} \times \np{0,0808} = \np{0,018471}$ soit \np{0,0185} à $10^{-4}$ près.
\item 
On a de même $p\left(\overline{N} \cap R\right) = p\left(N\right) \times p_{\overline{N}}(R) = \np{0,7714} \times \np{0,0127} = \np{0,009797}$ soit \np{0,0098} à $10^{-4}$ près.

D'après la loi des probabilités totales :

$p(R) = p(N \cap R) + p\left(\overline{N} \cap R\right) \approx \np{0,0185} + \np{0,0098}$

$p(R) \approx \np{0,0283}$.

\item 
On a $p_R(N) = \dfrac{p(R \cap N)}{p(R)} = \dfrac{p(N \cap R)}{p(R)} \approx \dfrac{\np{0,0185}}{\np{0,0283}} \approx \np{0,65371}$ soit \np{0,6537} à $10^{-4}$ près.
\end{enumerate}

\textbf{Partie B}

\medskip









\medskip

\begin{enumerate}
\item 

$X$ suit donc la loi binomiale $\mathcal{B}(n = 500, \: p = 0,65)$.
\item 
On calcule $\binom{500}{325} \times 0,65^{325} \times 0,35^{175} \approx \np{0,037384}$ soit \np{0,0374} à $10^{-4}$ près. (Utiliser la fonction binomiale de la calculatrice  si les capacités de calcul de celle-ci sont dépassées).
\item 
On a $p(X \geqslant 325) = 1 - p(X \leqslant 324)$ soit d'après la calculatrice \np{0,47944}, donc 

$p(X \geqslant 325) \approx 1 - \np{0,4794} = \np{0,5206}$ à $10^{-4}$ près.
\end{enumerate}

\textbf{Partie C}

\medskip







\medskip

\begin{enumerate}
\item 
On a $p_n = (1 - 0,65)^n = 0,35^n$.

\item On a $q_n = 1 - p_n$.

On cherche donc $n$ tel que $1 - p_n \geqslant \np{0,9999} \iff p_n \leqslant \np{0,0001}$, soit par croissance de la fonction logarithme népérien :

$n\ln 0,35 \leqslant \ln 0\np{0,0001} \iff n \geqslant \dfrac{\ln \np{0,0001}}{\ln 0,35}$ car $\dfrac{1}{\ln 0,35} < 0$.

D'après la calculatrice $\dfrac{\ln \np{0,0001}}{\ln 0,35} \approx 8,8$.

Il faut prendre au minimum $n = 9$.



\end{enumerate}

\end{document}
